\chapter{La route de Phandalin}

\begin{center}
    \textit{Mardi 25 août 2026}
\end{center}

\vspace{1em}

Waterdeep derrière eux, Kael, Morrÿs, Obélik et Alan prirent la route de Phandalin.

Ils n'étaient cependant pas tout à fait seuls.

Le capitaine des gardes, qui avait survécu tant bien que mal à sa rencontre avec la créature invisible, les accompagnait. Il n'était certes pas au meilleur de sa forme après avoir été transpercé à plusieurs reprises, mais il était vivant et suffisamment remis pour prendre la route avec eux.

Le voyage leur donna enfin l'occasion de faire un peu mieux connaissance.

Ils apprirent notamment que le capitaine avait un nom : Virgile.

La conversation revint naturellement sur les événements de la prison et sur cette immense faille dont personne n'avait pu apercevoir le fond. Virgile leur révéla alors qu'il avait déjà essayé d'en sonder la profondeur avec les moyens dont il disposait.

Il y avait jeté une chope.

Cette révélation provoqua une réaction immédiate d'Alan. En sa qualité de tavernier, sacrifier ainsi une chope dans un gouffre insondable lui paraissait manifestement beaucoup plus préoccupant que la méthode scientifique employée.

Virgile apporta toutefois une précision essentielle : avant de la jeter, il avait pris soin d'en vider intégralement le contenu.

On ne gâche pas.

La réponse sembla rendre l'expérience au moins légèrement plus acceptable.

\bigskip

Après quelque temps de marche, le groupe fit halte dans une taverne.

Le problème était que Kael, Morrÿs, Obélik et Alan portaient toujours les guenilles qui leur avaient servi de vêtements en prison. Pour quiconque prenait le temps de les observer, ils n'avaient donc pas exactement l'apparence de paisibles voyageurs en route vers Phandalin.

Virgile en était parfaitement conscient.

Avant d'entrer, il leur demanda de faire profil bas et, une fois à l'intérieur, de rester tranquillement à leur table. Lui se chargerait du reste.

Le capitaine alla commander pour tout le monde et prit soin de préciser au tavernier que ces gens étaient avec lui.

La présence d'un capitaine de la garde suffisait peut-être à éviter les questions directes.

Elle ne suffisait manifestement pas à empêcher la curiosité.

Autour d'eux, plusieurs clients observaient discrètement les quatres prisonniers en guenilles. Des regards se croisaient, suivis de quelques chuchotements à voix basse. Personne ne vint cependant leur demander directement ce que quatres individus manifestement sortis d'une cellule faisaient attablés dans une taverne en compagnie d'un officier de la garde.

Nos aventuriers appliquèrent donc, pour une fois, un plan particulièrement simple avec un certain succès :

rester assis et ne rien faire de suspect.

La pause leur permit surtout de poursuivre leurs discussions.

Alan leur en apprit davantage sur les circonstances assez particulières qui l'avaient conduit en prison. À Waterdeep, il avait régulièrement été confondu avec l'ancien Lord de la ville, \textit{[NOM DE L'ANCIEN LORD À COMPLÉTER]}.

Mais ce n'était pas cette ressemblance qui avait directement provoqué son arrestation.

Alan avait été emprisonné pour avoir posé des affiches dans les rues afin de faire la publicité de sa taverne.

Une affaire manifestement suffisamment grave pour justifier son séjour derrière les barreaux.

La fameuse taverne n'était d'ailleurs pas encore véritablement la sienne. Alan n'en était pour l'instant que le gérant. Il comptait néanmoins bien en devenir un jour le propriétaire et travaillait dans l'espoir de pouvoir finalement la racheter.

Virgile, de son côté, expliqua ce qu'il comptait faire une fois arrivé à Phandalin.

Ses blessures nécessitaient encore des soins et il connaissait justement une organisation susceptible de l'aider : les Faiseurs de Veuves.

Le nom n'était pas particulièrement rassurant.

Il s'agissait pourtant d'une guilde combattant pour le bien du royaume.

Son sinistre surnom ne venait pas de ses intentions, mais de l'échec de sa toute première mission. Chargés de sauver des otages, ses membres n'avaient pas réussi à tous les ramener. L'événement avait marqué les esprits et leur avait valu ce nom peu glorieux.

Ils avaient depuis accompli suffisamment d'actions pour démontrer qu'ils œuvraient réellement pour le royaume.

Ils n'avaient simplement jamais réussi à faire oublier leur première réputation.

Les Faiseurs de Veuves étaient restés les Faiseurs de Veuves.

\bigskip

La pause terminée, le groupe reprit la route de Phandalin.

Avant de partir, Alan accorda tout de même un dernier regard à la taverne dans laquelle ils venaient de manger.

Ce n'était pas trop mal.

Mais clairement moins bien que sa taverne.

Le voyage se poursuivit sans événement majeur jusqu'à ce que Phandalin apparaisse enfin au loin.

Quelque chose attira immédiatement leur attention.

De la fumée s'élevait au-dessus de la ville.

Obélik observa la colonne sombre et trouva immédiatement l'explication la plus vraisemblable :

\begin{quote}
    \itshape
    « Ils ont déjà allumé le barbecue. »
\end{quote}

L'hypothèse avait au moins le mérite de l'optimisme.

À mesure qu'ils approchaient, ils découvrirent également qu'ils n'arriveraient pas à Phandalin sans être accueillis.

Une elfe les attendait.

\illustration[0.3\textwidth]{Lyra_Sunstricke_pose.png}
{\textit{L'étrange elfe qui attendait les aventuriers aux abords de Phandalin.}}
{lyra_sunstricke}

Elle ne semblait ni surprise de les voir, ni particulièrement intriguée par leur arrivée.

Au contraire.

Elle leur annonça qu'elle les attendait.

Et ce, depuis un certain temps déjà.

\begin{narrateur}
À cet instant, une question aurait sans doute mérité d'être posée : comment une elfe rencontrée pour la première fois aux portes de Phandalin pouvait-elle savoir que quatre prisonniers évadés de Waterdeep, accompagnés d'un capitaine récemment transpercé à plusieurs reprises, allaient arriver ici ?

Nos aventuriers auraient tout le loisir de s'en inquiéter.
\end{narrateur}

% À compléter :
% - Nom de l'ancien Lord de Waterdeep avec lequel Alan est régulièrement confondu.
% - Date / durée exacte depuis laquelle l'elfe affirme attendre le groupe.